% Template voor een kampprogramma van de
% Stichting Vierkant voor Wiskunde. Copyright 2000-2005,2011-2012:
% M.P. Alberts, J.C.C. Langeveld, M. Hendriks, W.J. Palenstijn.
% Meest recente update: 1-6-2012

% De 'begeleider' optie produceert een begeleidersversie.
% De 'deelnemer' optie produceert een deelnemersversie (zonder antwoorden).
% Zonder 'begeleider' en 'deelnemer' wordt een voorlopige versie gemaakt,
%  met antwoorden, en met de huidige datum aangegeven.

% Er is ook een 'kampA'-optie die verschillende dingen doet. Zie README.txt .

%\documentclass[begeleider]{kampprogramma}
%\documentclass[deelnemer]{kampprogramma}
%\documentclass[deelnemer,kampA]{kampprogramma}
\documentclass{kampprogramma}

% ------- Voer hier de titel in
\titel{Template}
% -------

% ------- Voer hier de auteurs in
\auteurs{
Maurice Alberts\\
Maxim Hendriks\\
Willem Jan Palenstijn
}
% -------

% ------- Hier de naam van het (.eps) bestand met de kaftillustratie
% Je kan dit commando ook weglaten als je nog geen plaatje hebt.
\illustratie{vierkantlogo}
% -------

% ------- Hier het kamp waarvoor het programma is. (A/B/C + jaar)
\kamp{Zomerkamp C 2012}
% -------


\begin{document}

\voorpagina

\voorwoord
% ------- Schrijf hier een voorwoord
Hier het voorwoord van het kampprogramma
% -------

% ------- Legenda
%
% In de legenda worden symbolen voor opgaven uitgelegd.
% Alleen de daadwerkelijk in het programma gebruikte symbolen worden uitgelegd.
%
% De symbolen hebben de volgende betekenis:
% \ster : moeilijke opgave
% \vinger : deze opgave heeft een hint
% \schaar : dit is een praktische opdracht
% \gr : je mag/moet een grafische rekenmachine gebruiken voor deze opgave
% \discussie : dit is een discussie-opgave
%
% Zie het stukje 'Opgave' hieronder over hoe je de symbolen bij een opgave zet

\legenda


\inhoudsopgave

% ------- Hier begint het eigenlijke programma

\chapter{Het eerste hoofdstuk}
Hier de hoofdtekst van het kampprogramma. In het programma zullen opgaven
staan. Hieronder volgen uitleg en een paar voorbeelden.

%-------- Opgave
% Achter begin{opgave} kun je tussen vierkante haken [] symbolen
% uit de legenda toevoegen. In de eerste voorbeeldopgave worden
% alle symbolen meegegeven.
%
% Tussen dezelfde vierkante haken [] kun je met een getal het aantal
% regels opgeven dat na de opgave als antwoordruimte wordt
% geboden. Zie de tweede voorbeeldopgave.
%
% Wil je zowel legenda-symbolen als antwoordregels toevoegen, dan zet
% je binnen de vierkante haken [] tussen de symbolen en het getal een komma.
% Zie de derde voorbeeldopgave.
% --------

\begin{opgave}[\ster\vinger\schaar\gr]
Vul deze tekst aan tot een volledig kamprogramma
\begin{hint}
Dit is meer werk dan je zou denken!
\end{hint}
\begin{antwoord}
Geen antwoord voor deze opgave helaas\ldots
\end{antwoord}
\end{opgave}

Dit is wel een mooie plek voor een plaatje om wat ruimte te vullen:

\begin{center}
\includegraphics[width=8cm]{vierkantlogo}
\end{center}

\begin{opgave}[\ster]
Bewijs dat $a^3 + b^3 = c^3$ geen oplossingen heeft voor gehele $a$,
$b$ en $c$ ongelijk aan nul.
\begin{antwoord}
Dit is een direct gevolg van de laatste stelling van Fermat.
\end{antwoord}
\end{opgave}

\begin{opgave}
Het is ook mogelijk om een opgave in subopgaven te verdelen.
Antwoorden en hints kan je zowel in de subopgaven als in de hoofdopgave
zetten. Deze introductie is optioneel.
\begin{subopgave}
Dit is de eerste subopgave.
\begin{antwoord}
Geen antwoord nodig.
\end{antwoord}
\end{subopgave}
\begin{subopgave}[\schaar]
Er kunnen ook symbolen bij individuele subopgaven staan.
\end{subopgave}
\end{opgave}

\chapter{Het echte werk}

This chapter intentionally left blank.


% ------- Hier eindigt het eigenlijke programma

\hintsantwoorden

\end{document}
